% !TEX program = xelatex

\documentclass[12pt,a4paper]{article}

\usepackage{fontspec}
\usepackage{polyglossia}
\setmainlanguage{russian}
\setotherlanguage{english}
\setmainfont{DejaVu Serif}
\newfontfamily\cyrillicfonttt{DejaVu Sans Mono}
\usepackage{amsmath,amssymb,amsfonts,amsthm,mathtools}
\usepackage{geometry}
\usepackage{booktabs}
\usepackage{tabularx}
\usepackage{titlesec}
\usepackage{cite}
\usepackage{microtype}
\usepackage{xcolor}
\usepackage{tcolorbox}
\usepackage{hyperref}

\microtypesetup{expansion=false}
\setlength{\emergencystretch}{3em}

\titleformat{\section}{\large\bfseries}{\thesection}{1em}{}
\titleformat{\subsection}{\normalsize\bfseries}{\thesubsection}{1em}{}
\titleformat{\subsubsection}{\small\bfseries}{\thesubsubsection}{1em}{}

\geometry{
    left=25mm,
    right=20mm,
    top=25mm,
    bottom=25mm
}

\hypersetup{
    colorlinks=true,
    linkcolor=blue,
    citecolor=blue,
    urlcolor=blue
}

% Стилизованные блоки для справочника
\newtcolorbox{theorybox}[1]{
    colback=gray!6!white,
    colframe=gray!75!black,
    fonttitle=\bfseries\small,
    title=#1,
    arc=1.5mm,
    boxrule=0.8pt,
    top=2mm, bottom=2mm, left=3mm, right=3mm
}

\newtcolorbox{warningbox}[1]{
    colback=red!4!white,
    colframe=red!70!black,
    fonttitle=\bfseries\small,
    title=#1,
    arc=1.5mm,
    boxrule=0.8pt,
    top=2mm, bottom=2mm, left=3mm, right=3mm
}

\begin{document}

\section*{ГЛАВА 2. РАВНОВЕСНЫЕ СВОЙСТВА ПОЛИМЕРНЫХ ЦЕПЕЙ}

Конформация свободносочлененной цепи представляется набором из $(N + 1)$ векторов положений точек сочленения $\{\boldsymbol{R}_n\} = (\boldsymbol{R}_0 \dots \boldsymbol{R}_N)$, или по-другому~--- набором векторов связей $\{\boldsymbol{r}_n\} = (\boldsymbol{r}_1 \dots \boldsymbol{r}_N)$, где
\begin{equation}
\boldsymbol{r}_n = \boldsymbol{R}_n - \boldsymbol{R}_{n-1}, \quad n = 1, 2, \dots, N. \tag{2.1}
\end{equation}
Так как векторы связей $\boldsymbol{r}_n$ независимы друг от друга, функция распределения конформаций полимерной цепи запишется в виде
\begin{equation}
\Psi(\{\boldsymbol{r}_n\}) = \prod_{n=1}^N \psi(\boldsymbol{r}_n), \tag{2.2}
\end{equation}
где $\psi(\boldsymbol{r})$ обозначает случайное распределение вектора постоянной длины $b_0$:
\begin{equation}
\psi(\boldsymbol{r}) = \frac{1}{4\pi b_0^2}\delta(|\boldsymbol{r}| - b_0). \tag{2.3}
\end{equation}
Это распределение нормировано так, что
\begin{equation}
\int \!\mathrm{d}\boldsymbol{r}\,\psi(\boldsymbol{r}) = 1. \tag{2.4}
\end{equation}
Чтобы охарактеризовать размер полимерной цепи, рассмотрим вектор между концами цепи $\boldsymbol{R}$, где
\begin{equation}
\boldsymbol{R} = \boldsymbol{R}_N - \boldsymbol{R}_0 = \sum_{n=1}^N \boldsymbol{r}_n. \tag{2.5}
\end{equation}
Поскольку $\langle\boldsymbol{r}_n\rangle = 0$, то и $\langle\boldsymbol{R}\rangle = 0$, но $\langle\boldsymbol{R}^2\rangle$ имеет конечное значение, которое может быть использовано в качестве характеристики длины цепи. Определим $\overline{R}$ соотношением
\begin{equation}
\overline{R} \equiv \langle\boldsymbol{R}^2\rangle^{1/2} = \langle(\boldsymbol{R}_N - \boldsymbol{R}_0)^2\rangle^{1/2}. \tag{2.6}
\end{equation}
Используя уравнение (2.5), для $\langle\boldsymbol{R}^2\rangle$ получаем
\begin{equation}
\langle\boldsymbol{R}^2\rangle = \sum_{n,m=1}^N \langle\boldsymbol{r}_n \cdot \boldsymbol{r}_m\rangle = \sum_{n=1}^N \langle\boldsymbol{r}_n^2\rangle + 2\sum_{n>m}\langle\boldsymbol{r}_n \cdot \boldsymbol{r}_m\rangle = Nb_0^2, \tag{2.7}
\end{equation}
так как при $n \ne m$ $\langle\boldsymbol{r}_n \cdot \boldsymbol{r}_m\rangle = \langle\boldsymbol{r}_m\rangle \cdot \langle\boldsymbol{r}_n\rangle = 0$. Таким образом, $\overline{R}$ дается как
\begin{equation}
\overline{R} = \sqrt{N}b_0. \tag{2.8}
\end{equation}

\subsubsection*{2.1.2. Обобщенные модели случайных блужданий}

Хотя свободносочлененная цепь~--- очень простая модель, основной ее результат $\langle\boldsymbol{R}^2\rangle \propto N$ оказывается справедливым и для более общих моделей. Рассмотрим, например, показанную на рис.~2.2 модель\footnote{В литературе жесткость часто представляется длиной куновского статистического сегмента $b_K$, определяемой как $b_K = \langle\boldsymbol{R}^2\rangle / R_{\max}$, где $R_{\max}$ есть максимальная длина вектора между концами.} \dots

\subsubsection*{2.1.3. Распределение вектора между концами цепи}

В качестве следующего шага рассмотрим распределение вектора, соединяющего концы цепи, для модели случайных блужданий. Пусть $\Phi(\boldsymbol{R}, N)$ есть функция распределения вероятности того, что вектор между концами цепи, состоящей из $N$ связей, равен $\boldsymbol{R}$. Если распределение конформаций описывается функцией $\Psi(\{\boldsymbol{r}_n\})$, то $\Phi(\boldsymbol{R}, N)$ рассчитывается как
\begin{equation}
\Phi(\boldsymbol{R}, N) = \int\!\mathrm{d}\boldsymbol{r}_1\int\!\mathrm{d}\boldsymbol{r}_2 \dots \int\!\mathrm{d}\boldsymbol{r}_N \, \delta\left(\boldsymbol{R} - \sum_{n=1}^N \boldsymbol{r}_n\right)\Psi(\{\boldsymbol{r}_n\}), \tag{2.18}
\end{equation}
что можно переписать, используя тождество
\begin{equation}
\delta(\boldsymbol{r}) = \frac{1}{(2\pi)^3}\int\!\mathrm{d}\boldsymbol{k}\,\mathrm{e}^{i\boldsymbol{k}\cdot\boldsymbol{r}}, \tag{2.19}
\end{equation}
в виде
\begin{equation}
\Phi(\boldsymbol{R}, N) = \frac{1}{(2\pi)^3}\int\!\mathrm{d}\boldsymbol{k}\int\!\mathrm{d}\boldsymbol{r}_1\int\!\mathrm{d}\boldsymbol{r}_2 \dots \int\!\mathrm{d}\boldsymbol{r}_N \times \exp\left(ik\cdot\left(\boldsymbol{R} - \sum_{n=1}^N \boldsymbol{r}_n\right)\right)\Psi(\{\boldsymbol{r}_n\}). \tag{2.20}
\end{equation}
Для свободносочлененной цепи уравнения (2.2) и (2.20) дают
\begin{align}
\Phi(\boldsymbol{R}, N) &= \frac{1}{(2\pi)^3}\int\!\mathrm{d}\boldsymbol{k}\,\mathrm{e}^{i\boldsymbol{k}\cdot\boldsymbol{R}}\int\!\mathrm{d}\boldsymbol{r}_1 \dots \mathrm{d}\boldsymbol{r}_N \prod_{n=1}^N \exp(-i\boldsymbol{k}\cdot\boldsymbol{r}_n)\psi(\boldsymbol{r}_n) = \nonumber\\
&= \frac{1}{(2\pi)^3}\int\!\mathrm{d}\boldsymbol{k}\,\mathrm{e}^{i\boldsymbol{k}\cdot\boldsymbol{R}}\left[\int\!\mathrm{d}\boldsymbol{r}\exp(-i\boldsymbol{k}\cdot\boldsymbol{r})\psi(\boldsymbol{r})\right]^N. \tag{2.21}
\end{align}
Интеграл по $\boldsymbol{r}$ оценивают, вводя полярные координаты $(r, \theta, \phi)$; причем ось отсчета для $\theta$ выбирают вдоль вектора $\boldsymbol{k}$, что дает
\begin{align}
\int\!\mathrm{d}\boldsymbol{r}\exp(-i\boldsymbol{k}\cdot\boldsymbol{r})\psi(\boldsymbol{r}) &= \frac{1}{4\pi b^2}\int_0^\infty\!\mathrm{d}r\,r^2\int_0^{2\pi}\!\mathrm{d}\phi\int_0^\pi\!\mathrm{d}\theta\sin\theta\exp(-ikr\cos\theta)\delta(r-b) = \nonumber\\
&= \frac{\sin kb}{kb}, \tag{2.22}
\end{align}
где $k = |\boldsymbol{k}|$, а $b_0$ заменяется на $b$ (так как $b_0 = b$ для свободносочлененной цепи). Из уравнений (2.21) и (2.22) имеем
\begin{equation}
\Phi(\boldsymbol{R}, N) = \frac{1}{(2\pi)^3}\int\!\mathrm{d}\boldsymbol{k}\,\mathrm{e}^{i\boldsymbol{k}\cdot\boldsymbol{R}}\left(\frac{\sin kb}{kb}\right)^N. \tag{2.23}
\end{equation}

Если $N$ велико, $((\sin kb)/kb)^N$ становится очень малым при условии, что $kb$ мало. При $kb \ll 1$ $((\sin kb)/kb)^N$ можно аппроксимировать в виде
\begin{equation}
\left(\frac{\sin kb}{kb}\right)^N = \left(1 - \frac{k^2 b^2}{6}\right)^N = \exp\left(-\frac{Nk^2 b^2}{6}\right). \tag{2.24}
\end{equation}
Это приближение выполняется также и при $kb \gg 1$, так как в этом случае обе части уравнения (2.24) приблизительно равны нулю. Таким образом, $\Phi(\boldsymbol{R}, N)$ рассчитывается как
\begin{equation}
\Phi(\boldsymbol{R}, N) = \frac{1}{(2\pi)^3}\int\!\mathrm{d}\boldsymbol{k}\,\mathrm{e}^{i\boldsymbol{k}\cdot\boldsymbol{R}}\exp\left(-\frac{Nk^2 b^2}{6}\right). \tag{2.25}
\end{equation}
Интеграл по $\boldsymbol{k}$ является стандартным гауссовым интегралом. (Сводка некоторых полезных гауссовых интегралов приведена в приложении~2.I.) Если $k_\alpha$ и $R_\alpha$ ($\alpha = x, y, z$) обозначают компоненты векторов $\boldsymbol{k}$ и $\boldsymbol{R}$, то, используя уравнение (2.I.2), имеем
\begin{align}
\Phi(\boldsymbol{R}, N) &= (2\pi)^{-3}\prod_{\alpha=x,y,z}\left[\int_{-\infty}^\infty\!\mathrm{d}k_\alpha\exp(ik_\alpha R_\alpha - Nk_\alpha^2 b^2/6)\right] = \nonumber\\
&= (2\pi)^{-3}\prod_{\alpha=x,y,z}\left(\frac{6\pi}{Nb^2}\right)^{1/2}\exp\left(-\frac{3}{2Nb^2}R_\alpha^2\right) = \nonumber\\
&= (3/2\pi Nb^2)^{3/2}\exp\left(-\frac{3\boldsymbol{R}^2}{2Nb^2}\right). \tag{2.26}
\end{align}
Таким образом, функция распределения вектора между концами является гауссовой.

Распределение (2.26) имеет ту искусственно возникшую особенность, что $|\boldsymbol{R}|$ может быть больше максимальной длины вытянутой цепи $Nb$. В литературе [1, 2, 4] можно найти более реалистичную функцию распределения, которая не имеет этой особенности. В этой книге, однако, такие сильно вытянутые состояния полимерных цепей не рассматриваются, и выражения (2.26) для наших целей достаточно.

Хотя вышеприведенный вывод проделан для свободносочлененной цепи, полученный результат справедлив и в более общем случае. Вообще можно показать, что, когда распределение конформаций описывается формулой (2.15), распределение вектора между концами $\boldsymbol{R}$ длинной цепи ($N \gg 1$) дается выражением (2.26).

\begin{center}
\textit{Рис. 2.3. Цепь, разделенная на $\bar{N}$ субмолекул.}
\end{center}

Этот результат следует из центральной предельной теоремы в математической статистике [5].

Для доказательства этого предположим, что цепь разделена на $\bar{N}$ субмолекул, каждая из которых содержит $\lambda$ связей (рис.~2.3). Ясно, что
\begin{equation}
\bar{N} = N/\lambda. \tag{2.27}
\end{equation}
Пусть $\bar{\boldsymbol{r}}_n$ есть вектор между концами $n$-й субмолекулы. Вектор между концами всей цепи запишется как
\begin{equation}
\boldsymbol{R} = \sum_{n=1}^{\bar{N}}\bar{\boldsymbol{r}}_n. \tag{2.28}
\end{equation}
Теперь если $N$ очень велико, то и $\bar{N}$, и $\lambda$ могут быть взяты достаточно большими, так что $\bar{N} \gg 1$ и $\lambda \gg 1$. Если $\lambda \gg 1$, то векторы $\bar{\boldsymbol{r}}_n$ становятся независимыми друг от друга, и распределение $\{\bar{\boldsymbol{r}}_n\}$ можно записать как
\begin{equation}
\Psi(\{\bar{\boldsymbol{r}}_n\}) = \prod_{n=1}^{\bar{N}}\bar{\psi}(\bar{\boldsymbol{r}}_n). \tag{2.29}
\end{equation}
При выводе распределения $\boldsymbol{R}$ из уравнений (2.28) и (2.29) не требуется знания действительного вида функции $\bar{\psi}(\bar{\boldsymbol{r}})$. Вновь используя тождества (2.21) и (2.29), имеем
\begin{equation}
\Phi(\boldsymbol{R}, N) = \frac{1}{(2\pi)^3}\int\!\mathrm{d}\boldsymbol{k}\,\mathrm{e}^{i\boldsymbol{k}\cdot\boldsymbol{R}}A(k)^{\bar{N}}, \tag{2.30}
\end{equation}
где
\begin{equation}
A(k) = \int\!\mathrm{d}\boldsymbol{r}\exp(-i\boldsymbol{k}\cdot\boldsymbol{r})\bar{\psi}(\boldsymbol{r}). \tag{2.31}
\end{equation}
Так как $\bar{\psi}(\boldsymbol{r})$ зависит только от $|\boldsymbol{r}|$, интегрирование по направлению $\boldsymbol{r}$ может быть выполнено:
\begin{align}
A(k) &= \int_0^\infty\!\mathrm{d}r\,r^2\int_0^\pi\!\mathrm{d}\theta\sin\theta\int_0^{2\pi}\!\mathrm{d}\phi\exp(-ikr\cos\theta)\bar{\psi}(r) = \nonumber\\
&= \int_0^\infty\!\mathrm{d}r\,4\pi r^2\frac{\sin kr}{kr}\bar{\psi}(r) = \left\langle\frac{\sin kr}{kr}\right\rangle_{\!\bar{\psi}}, \tag{2.32}
\end{align}
где
\begin{equation}
\langle\dots\rangle_{\bar{\psi}} \equiv \int_0^\infty\!\mathrm{d}r\,4\pi r^2\bar{\psi}(r)\dots \tag{2.33}
\end{equation}
Поскольку нас интересует область малых $k$, мы можем оценить $A(k)$, разлагая $(\sin kr)/kr$ в ряд по $k$:
\begin{equation}
A(k) = \left\langle 1 - \frac{1}{6}k^2 r^2 + \dots\right\rangle_{\!\bar{\psi}} = 1 - \frac{1}{6}k^2\langle r^2\rangle_{\bar{\psi}}. \tag{2.34}
\end{equation}
При $\lambda \gg 1$ $\langle r^2\rangle_{\bar{\psi}}$ дается выражением $\lambda b^2$, откуда
\begin{equation}
A(k) = 1 - \frac{1}{6}\lambda b^2 k^2. \tag{2.35}
\end{equation}
Тогда $\Phi(\boldsymbol{R}, N)$ оценивается как
\begin{align}
\Phi(\boldsymbol{R}, N) &= \frac{1}{(2\pi)^3}\int\!\mathrm{d}\boldsymbol{k}\,\mathrm{e}^{i\boldsymbol{k}\cdot\boldsymbol{R}}\left(1 - \frac{1}{6}\lambda b^2 k^2\right)^{\bar{N}} = \nonumber\\
&= \frac{1}{(2\pi)^3}\int\!\mathrm{d}\boldsymbol{k}\,\mathrm{e}^{i\boldsymbol{k}\cdot\boldsymbol{R}}\exp\left(-\frac{1}{6}\lambda\bar{N}b^2 k^2\right) = \nonumber\\
&= \left(\frac{3}{2\pi\lambda\bar{N}b^2}\right)^{3/2}\exp\left(-\frac{3\boldsymbol{R}^2}{2\lambda\bar{N}b^2}\right). \tag{2.36}
\end{align}
Поскольку $\lambda\bar{N} = N$, формула (2.36) полностью согласуется с выражением (2.26).

\subsection*{2.2. ГАУССОВА ЦЕПЬ}

Мы убедились, что в статистическом распределении вектора между концами локальная структура цепи проявляется только через эффективную длину связи $b$. В общем случае справедливо утверждение: локальная структура цепи определяет только эффективную длину последней и никаким другим путем в рассматриваемой проблеме не проявляется. Поэтому если мы интересуемся общими свойствами полимеров, мы можем исходить из простейшей допустимой модели.

Рассмотрим цепь, в которой длина связи имеет гауссово распределение:
\begin{equation}
\psi(\boldsymbol{r}) = \left[\frac{3}{2\pi b^2}\right]^{3/2}\exp\left(-\frac{3\boldsymbol{r}^2}{2b^2}\right), \tag{2.37}
\end{equation}
так что
\begin{equation}
\langle\boldsymbol{r}^2\rangle = b^2. \tag{2.38}
\end{equation}
Функция распределения конформаций такой цепи дается выражением
\begin{align}
\Psi(\{\boldsymbol{r}_n\}) &= \prod_{n=1}^N\left[\frac{3}{2\pi b^2}\right]^{3/2}\exp\left[-\frac{3\boldsymbol{r}_n^2}{2b^2}\right] = \nonumber\\
&= \left[\frac{3}{2\pi b^2}\right]^{3N/2}\exp\left[-\sum_{n=1}^N\frac{3(\boldsymbol{R}_n - \boldsymbol{R}_{n-1})^2}{2b^2}\right]. \tag{2.39}
\end{align}
Подобная цепь называется гауссовой. Гауссова цепь не описывает строго локальную структуру полимера, но правильно воспроизводит его свойства на больших масштабах длин. Основное преимущество гауссовой цепи как модели заключается в том, что математически с нею оперировать гораздо легче, чем с любой из моделей, рассмотренных в разд.~2.1.

Гауссова цепь часто представляется механической моделью (рис.~2.4); принимается, что $(N + 1)$ «бусинок» связаны гармонической пружиной, потенциальная энергия которой дается выражением
\begin{equation}
U_0(\{\boldsymbol{R}_n\}) = \frac{3}{2b^2}k_B T \sum_{n=1}^N(\boldsymbol{R}_n - \boldsymbol{R}_{n-1})^2. \tag{2.40}
\end{equation}
В равновесии больцмановское распределение в рамках такой модели в точности то же самое, что и в формуле (2.39).

Важным свойством гауссовой цепи является то, что распределение вектора $\boldsymbol{R}_n - \boldsymbol{R}_m$ между любыми двумя единицами $n$ и $m$ является гауссовым и дается выражением
\begin{equation}
\Phi(\boldsymbol{R}_n - \boldsymbol{R}_m,\, n - m) = \left[\frac{3}{2\pi b^2 |n - m|}\right]^{3/2}\exp\left[-\frac{3(\boldsymbol{R}_n - \boldsymbol{R}_m)^2}{2|n - m|b^2}\right]. \tag{2.41}
\end{equation}
Это следует из свойств гауссова интеграла (см. приложение~2.I). Соответственно при любых $n$ и $m$
\begin{equation}
\langle(\boldsymbol{R}_n - \boldsymbol{R}_m)^2\rangle = |n - m|b^2. \tag{2.42}
\end{equation}
Индекс $n$ для гауссовой цепи часто рассматривается как непрерывная переменная. В этом случае $\boldsymbol{R}_n - \boldsymbol{R}_{n-1}$ заменяется на $\partial\boldsymbol{R}_n/\partial n$, и выражение (2.39) записывается как
\begin{equation}
\Psi[\boldsymbol{R}_n] = \text{const}\,\exp\left[-\frac{3}{2b^2}\int_0^N\!\mathrm{d}n\left(\frac{\partial\boldsymbol{R}_n}{\partial n}\right)^2\right]. \tag{2.43}
\end{equation}
Это распределение известно как распределение Винера.

Математически существует некая тонкость при переходе от дискретного $n$ к непрерывному, но для наших целей достаточно понимать, что уравнение (2.43) является лишь иной формальной записью выражения (2.39). В этой книге дискретное и непрерывное $n$ рассматриваются как взаимозаменяемые. Правила перехода от дискретных переменных к непрерывным суммированы в табл.~2.1.

\begin{table}[h!]
\centering
\caption*{Таблица 2.1}
\begin{tabular}{ccc}
\toprule
Дискретное $n$ & & Непрерывное $n$ \\
\midrule
$\displaystyle\sum_{n=k_1}^{k_2}$ & $\to$ & $\displaystyle\int_{k_1}^{k_2}\!\mathrm{d}n$ \\
$\boldsymbol{R}_n - \boldsymbol{R}_{n-1}$ & $\to$ & $\partial\boldsymbol{R}_n/\partial n$ \\
$\boldsymbol{R}_{n+1} + \boldsymbol{R}_{n-1} - 2\boldsymbol{R}_n$ & $\to$ & $\partial^2\boldsymbol{R}_n/\partial n^2$ \\
$\delta_{nm}$ Дельта Кронекера & $\to$ & $\delta(n - m)$ Дельта-функция Дирака \\
$\displaystyle\frac{\partial}{\partial\boldsymbol{R}_n}$ & $\to$ & $\displaystyle\frac{\delta}{\delta\boldsymbol{R}_n}{}^{*)}$ \\
$\displaystyle\int\prod\mathrm{d}\boldsymbol{R}_n$ & $\to$ & $\displaystyle\int\!\delta\boldsymbol{R}_n{}^{**)}$ \\
\bottomrule
\end{tabular}
\end{table}
\noindent\footnotesize{${}^{*), **)}$ Эти символы относятся к функциональным производной и интегралу соответственно. В принятых обозначениях они записываются как $\delta/\delta\boldsymbol{R}(n)$ и $\int\!\delta\boldsymbol{R}(n)$, чтобы подчеркнуть, что $\boldsymbol{R}$ есть функция непрерывной переменной $n$. Однако, поскольку здесь дискретное и непрерывное представления взаимозаменяемы, мы используем символ $\boldsymbol{R}_n$ в обоих представлениях.}

\normalsize

\subsection*{2.3. КОНФОРМАЦИЯ ЦЕПИ ВО ВНЕШНЕМ ПОЛЕ}

\subsubsection*{2.3.1. Функция Грина}

Если имеется внешнее поле $U_e(\boldsymbol{r})$, действующее на каждый сегмент, то равновесное распределение гауссовой цепи видоизменяется с учетом множителя Больцмана:
\begin{equation}
\exp\left[-\frac{1}{k_B T}\int_0^N\!\mathrm{d}n\,U_e(\boldsymbol{R}_n)\right], \tag{2.44}
\end{equation}
и конформационная функция распределения приобретает вид
\begin{equation}
\Psi[\boldsymbol{R}_n] \propto \exp\left[-\frac{3}{2b^2}\int_0^N\!\mathrm{d}n\left(\frac{\partial\boldsymbol{R}_n}{\partial n}\right)^2 - \frac{1}{k_B T}\int_0^N\!\mathrm{d}n\,U_e(\boldsymbol{R}_n)\right]. \tag{2.45}
\end{equation}
При рассмотрении статистических свойств такой системы удобно ввести функцию Грина, определяемую соотношением [6, 7]:
\begin{equation}
G(\boldsymbol{R}, \boldsymbol{R}'; N) \equiv \frac{\displaystyle\int_{\boldsymbol{R}_0 = \boldsymbol{R}'}^{\boldsymbol{R}_N = \boldsymbol{R}}\!\delta\boldsymbol{R}_n\exp\left[-\frac{3}{2b^2}\int_0^N\!\mathrm{d}n\left(\frac{\partial\boldsymbol{R}_n}{\partial n}\right)^2 - \frac{1}{k_B T}\int_0^N\!\mathrm{d}n\,U_e(\boldsymbol{R}_n)\right]}{\displaystyle\int\!\mathrm{d}\boldsymbol{R}\int_{\boldsymbol{R}_0 = \boldsymbol{R}'}^{\boldsymbol{R}_N = \boldsymbol{R}}\!\delta\boldsymbol{R}_n\exp\left[-\frac{3}{2b^2}\int_0^N\!\mathrm{d}n\left(\frac{\partial\boldsymbol{R}_n}{\partial n}\right)^2\right]}, \tag{2.46}
\end{equation}
где $\boldsymbol{R}_N = \boldsymbol{R}$ и $\boldsymbol{R}_0 = \boldsymbol{R}'$ в функциональном интеграле означает, что интеграл берется по всем конформациям, которые начинаются в $\boldsymbol{R}'$ и заканчиваются в $\boldsymbol{R}$. (Заметим, что знаменатель в уравнении (2.46) не зависит от $\boldsymbol{R}$ и $\boldsymbol{R}'$.)

При $U_e = 0$ $G(\boldsymbol{R}, \boldsymbol{R}'; N)$ сводится к гауссовой функции распределения
\begin{equation}
G(\boldsymbol{R} - \boldsymbol{R}'; N) = \left(\frac{2\pi N b^2}{3}\right)^{-3/2}\exp\left(-\frac{3(\boldsymbol{R} - \boldsymbol{R}')^2}{2Nb^2}\right). \tag{2.47}
\end{equation}

\begin{center}
\textit{Рис. 2.5. Цепь, которая начинается в $\boldsymbol{R}'$, проходит через $\boldsymbol{R}''$ за $n$ шагов и заканчивается в $\boldsymbol{R}$ за $N$ шагов.}
\end{center}

В общем случае при $U_e \ne 0$ $G(\boldsymbol{R}, \boldsymbol{R}'; N)$ представляет статистический вес (или статистическую сумму) цепи, которая выходит из $\boldsymbol{R}'$ и заканчивается в $\boldsymbol{R}$ за $N$ шагов. Статистический вес для всех возможных конформаций дается выражением
\begin{equation}
Z = \int\!\mathrm{d}\boldsymbol{R}\,\mathrm{d}\boldsymbol{R}'\,G(\boldsymbol{R}, \boldsymbol{R}'; N). \tag{2.48}
\end{equation}
Из определения $G(\boldsymbol{R}, \boldsymbol{R}'; N)$ вытекает справедливость следующего тождества:
\begin{equation}
G(\boldsymbol{R}, \boldsymbol{R}'; N) = \int\!\mathrm{d}\boldsymbol{R}''\,G(\boldsymbol{R}, \boldsymbol{R}''; N - n)G(\boldsymbol{R}'', \boldsymbol{R}'; n) \quad (\text{при } 0 < n < N). \tag{2.49}
\end{equation}
Физический смысл этого соотношения ясен. Множитель $G(\boldsymbol{R}, \boldsymbol{R}''; N - n)G(\boldsymbol{R}'', \boldsymbol{R}'; n)$ представляет собой статистический вес цепи, которая начинается в $\boldsymbol{R}'$, проходит через $\boldsymbol{R}''$ за $n$ шагов и заканчивается в $\boldsymbol{R}$ за $N$ шагов (рис.~2.5). Интегрирование этого статистического веса по всем $\boldsymbol{R}''$ дает статистический вес цепи, которая начинается в $\boldsymbol{R}'$ и заканчивается в $\boldsymbol{R}$.

Если $G(\boldsymbol{R}, \boldsymbol{R}'; N)$ известна, то легко рассчитать среднее значение любой физической величины $A$. Если $A$ зависит только от положения $n$-го сегмента, то
\begin{equation}
\langle A(\boldsymbol{R}_n)\rangle = \frac{\displaystyle\int\!\mathrm{d}\boldsymbol{R}_N\mathrm{d}\boldsymbol{R}_n\mathrm{d}\boldsymbol{R}_0\,G(\boldsymbol{R}_N, \boldsymbol{R}_n; N - n)G(\boldsymbol{R}_n, \boldsymbol{R}_0; n)A(\boldsymbol{R}_n)}{\displaystyle\int\!\mathrm{d}\boldsymbol{R}_N\mathrm{d}\boldsymbol{R}_0\,G(\boldsymbol{R}_N, \boldsymbol{R}_0; N)}. \tag{2.50}
\end{equation}
Сходным образом если $A$ зависит от $\boldsymbol{R}_n$ и $\boldsymbol{R}_m$ (принимается, что $n > m$), то
\begin{align}
\langle A(\boldsymbol{R}_n, \boldsymbol{R}_m)\rangle &= \int\!\mathrm{d}\boldsymbol{R}_N\mathrm{d}\boldsymbol{R}_n\mathrm{d}\boldsymbol{R}_m\mathrm{d}\boldsymbol{R}_0\,G(\boldsymbol{R}_N, \boldsymbol{R}_n; N - n)G(\boldsymbol{R}_n, \boldsymbol{R}_m; n - m) \nonumber\\
&\quad \times G(\boldsymbol{R}_m, \boldsymbol{R}_0; m)A(\boldsymbol{R}_n, \boldsymbol{R}_m)\times\left[\int\!\mathrm{d}\boldsymbol{R}_N\mathrm{d}\boldsymbol{R}_0\,G(\boldsymbol{R}_N, \boldsymbol{R}_0; N)\right]^{-1}. \tag{2.51}
\end{align}

Хотя функция Грина $G(\boldsymbol{R}, \boldsymbol{R}'; N)$ имеет физический смысл только для $N > 0$, удобно определить $G(\boldsymbol{R}, \boldsymbol{R}'; N)$ при $N < 0$ таким образом, что
\begin{equation}
G(\boldsymbol{R}, \boldsymbol{R}'; N) = 0 \quad \text{при } N < 0. \tag{2.52}
\end{equation}
При таком определении $G(\boldsymbol{R}, \boldsymbol{R}'; N)$ удовлетворяет простому дифференциальному уравнению
\begin{equation}
\left(\frac{\partial}{\partial N} - \frac{b^2}{6}\frac{\partial^2}{\partial\boldsymbol{R}^2} + \frac{1}{k_B T}U_e(\boldsymbol{R})\right)G(\boldsymbol{R}, \boldsymbol{R}'; N) = \delta(\boldsymbol{R} - \boldsymbol{R}')\delta(N). \tag{2.53}
\end{equation}
Вывод его дан в приложении~2.II. Произведение дельта-функций $\delta(\boldsymbol{R} - \boldsymbol{R}')\delta(N)$ в правой части учитывает граничные условия
\begin{equation}
G(\boldsymbol{R}, \boldsymbol{R}'; N) = 0 \quad \text{при } N < 0 \quad \text{и} \quad G(\boldsymbol{R}, \boldsymbol{R}'; N = 0) = \delta(\boldsymbol{R} - \boldsymbol{R}'). \tag{2.54}
\end{equation}

\subsubsection*{2.3.2. Пример: цепь в ящике}

Как пример применения функции Грина рассмотрим полимерную цепь, заключенную в ящик объемом $V = L_x L_y L_z$ (рис.~2.6) [8]. Ограничение на конформацию цепи со стороны стенок описывается посредством внешнего потенциала $U_e$, который имеет бесконечную величину вне ящика и равен нулю внутри ящика. Другой способ состоит в описании с помощью граничного условия:
\begin{equation}
G(\boldsymbol{R}, \boldsymbol{R}'; N) = 0, \quad \text{если $\boldsymbol{R}$ находится на границе ящика.} \tag{2.55}
\end{equation}
Решение уравнения
\begin{equation}
\left(\frac{\partial}{\partial N} - \frac{b^2}{6}\frac{\partial^2}{\partial\boldsymbol{R}^2}\right)G(\boldsymbol{R}, \boldsymbol{R}'; N) = \delta(\boldsymbol{R} - \boldsymbol{R}')\delta(N) \tag{2.56}
\end{equation}
при граничном условии (2.55) находится стандартным методом [8]:
\begin{equation}
G(\boldsymbol{R}, \boldsymbol{R}'; N) = g_x(R_x, R_x'; N)g_y(R_y, R_y'; N)g_z(R_z, R_z'; N), \tag{2.57}
\end{equation}
где
\begin{equation}
g_x(R_x, R_x'; N) = \frac{2}{L_x}\sum_{p=1}^\infty \sin\left(\frac{p\pi R_x}{L_x}\right)\sin\left(\frac{p\pi R_x'}{L_x}\right)\exp\left(-\frac{p^2\pi^2 Nb^2}{6L_x^2}\right). \tag{2.58}
\end{equation}
Статистический вес тогда дается выражением
\begin{equation}
Z = \int\!\mathrm{d}\boldsymbol{R}\,\mathrm{d}\boldsymbol{R}'\,G(\boldsymbol{R}, \boldsymbol{R}'; N) = Z_x Z_y Z_z, \tag{2.59}
\end{equation}
где
\begin{align}
Z_x &= \int_0^{L_x}\!\mathrm{d}R_x\int_0^{L_x}\!\mathrm{d}R_x'\,g_x(R_x, R_x'; N) = \nonumber\\
&= \frac{8}{\pi^2}L_x\sum_{p=1,3,\dots}^\infty \frac{1}{p^2}\exp\left[-\frac{\pi^2 Nb^2 p^2}{6L_x^2}\right]. \tag{2.60}
\end{align}
Свободная энергия системы рассчитывается по соотношению
\begin{equation}
A = -k_B T \ln Z, \tag{2.61}
\end{equation}
так что давление на стенку, перпендикулярную оси $x$, равно
\begin{equation}
P_x = -\frac{1}{L_y L_z}\frac{\partial A}{\partial L_x}. \tag{2.62}
\end{equation}
Рассмотрим два предельных случая:

1) В случае когда полимерная цепь много меньше ящика (т.~е. $\sqrt{N}b \ll L_x, L_y, L_z$), $Z_x$ приближенно равно
\begin{equation}
Z_x \approx \frac{8}{\pi^2}L_x\sum_{p=1,3,\dots}\frac{1}{p^2} = L_x, \tag{2.63}
\end{equation}
откуда
\begin{equation}
P_x = \frac{k_B T}{L_x L_y L_z} = \frac{k_B T}{V} \quad \text{и аналогично} \quad P_y = P_z = \frac{k_B T}{V}, \tag{2.64}
\end{equation}
что просто является уравнением состояния идеального газа.

2) В случае когда размеры полимерного клубка много больше размеров ящика (т.~е. $\sqrt{N}b \gg L_x, L_y, L_z$), $Z_x$ по преимуществу определяется первым слагаемым в сумме и
\begin{equation}
Z_x \approx \frac{8}{\pi^2}L_x\exp\left[-\frac{\pi^2 Nb^2}{6L_x^2}\right], \tag{2.65}
\end{equation}
что дает
\begin{equation}
P_x = -\frac{1}{L_y L_z}\frac{\partial A}{\partial L_x} = \frac{k_B T}{V}\left(1 + \frac{\pi^2 Nb^2}{3L_x^2}\right) \approx \frac{\pi^2 Nb^2}{3L_x^2}\frac{k_B T}{V}. \tag{2.66}
\end{equation}
Поскольку в этом случае $P_x$ не равно $P_y$ и $P_z$ (если только не выполнено условие $L_x = L_y = L_z$), то сила, действующая на стенки, не является изотропной. Позднее будет показано, что эта анизотропия давления, обусловленная анизотропией ограничений на цепь, ответственна за целый ряд механических свойств, специфичных для полимеров.

Метод функции Грина применяется к различным задачам, связанным с полимерными цепями, которые находятся в неоднородных ситуациях (адсорбция полимерных молекул, цепи вблизи границы раздела фаз, полимерные молекулы в пористой среде и т.~д.). Много таких примеров можно найти в литературе [9, 10].

\subsection*{2.5.2. МОДЕЛЬ ЦЕПИ С УЧЕТОМ ИСКЛЮЧЕННОГО ОБЪЕМА}

В реальных полимерных цепях природа дальнего взаимодействия весьма сложна: оно включает стерические эффекты, силы притяжения Ван-дер-Ваальса, а также другие специфические взаимодействия, включающие молекулы растворителя. Однако, если речь идет о крупномасштабных свойствах, детали взаимодействия не существенны, важно лишь то, что эффект исключенного объема контролируется взаимодействием между удаленными частями цепи. Таким образом, взаимодействие между полимерными сегментами с номерами $n$ и $m$ может быть выражено функцией близкодействия в пространстве
\begin{equation}
k_B T \tilde{v}(\boldsymbol{R}_n - \boldsymbol{R}_m), \tag{2.85}
\end{equation}
которая в дальнейшем будет аппроксимироваться дельта-функцией [19]:
\begin{equation}
v k_B T\,\delta(\boldsymbol{R}_n - \boldsymbol{R}_m), \tag{2.86}
\end{equation}
где $v$ есть исключенный объем и имеет размерность объема\footnote{Читателя может беспокоить использование дельта-функции, которая расходится, когда $\boldsymbol{R}_n = \boldsymbol{R}_m$. Однако, как оказывается, все физически измеримые величины включают интегралы от такого рода потенциалов и обладают хорошими свойствами. (Абсолютная энтропия цепи, скажем, расходится, но разность энтропий, с которой имеют дело экспериментально, определяется вполне корректно.) Эта ситуация хорошо известна в квантовых теориях поля и подробно обсуждается в учебниках. Можно заверить читателя, что в этом пункте реальных трудностей не возникает.}.

Тогда полная энергия взаимодействия запишется как
\begin{equation}
U_1 = \frac{1}{2}v k_B T\int_0^N\!\mathrm{d}n\int_0^N\!\mathrm{d}m\,\delta(\boldsymbol{R}_n - \boldsymbol{R}_m). \tag{2.87}
\end{equation}
Используя локальную концентрацию сегментов
\begin{equation}
c(\boldsymbol{r}) = \sum_n \delta(\boldsymbol{r} - \boldsymbol{R}_n) = \int_0^N\!\mathrm{d}n\,\delta(\boldsymbol{r} - \boldsymbol{R}_n), \tag{2.88}
\end{equation}
формулу (2.87) можно переписать в виде
\begin{equation}
U_1 = \int\!\mathrm{d}\boldsymbol{r}\,\frac{1}{2}v k_B T c(\boldsymbol{r})^2. \tag{2.89}
\end{equation}
Это выражение показывает, что (2.87) является первым членом вириального разложения свободной энергии по локальной концентрации $c(\boldsymbol{r})$. Поэтому параметр исключенного объема $v$ можно рассматривать как вириальный коэффициент взаимодействия между сегментами.

В принципе это вириальное разложение можно продолжить, чтобы включить члены более высокого порядка, такие, как
\begin{equation}
U_1 = \int\!\mathrm{d}\boldsymbol{r}\left[\frac{1}{2}v k_B T c(\boldsymbol{r})^2 + \frac{1}{6}w k_B T c(\boldsymbol{r})^3 + \dots\right]. \tag{2.90}
\end{equation}
Однако членами более высокого порядка можно пренебречь, так как плотность сегментов внутри полимерного клубка мала: плотность сегментов оценивается как
\begin{equation}
\bar{c} \approx \frac{N}{R^3} \propto N^{1 - 3\nu} = N^{-4/5} \quad (\text{при } \nu = 3/5) \tag{2.91}
\end{equation}
и становится очень малой при больших $N$. Поэтому все существенные черты эффекта исключенного объема могут быть исследованы с помощью потенциала, даваемого выражением (2.87).

Для выбранной комбинации полимера и растворителя $v$ меняется с температурой и может стать равным нулю при некоторой температуре $\Theta$, называемой также температурой Флори. При температуре $\Theta$ цепь становится близкой к идеальной\footnote{Даже при $\Theta$-температуре цепь не является идеальной, поскольку имеется член, соответствующий трехчастичным столкновениям [20, 21]. Однако эффект трехчастичных столкновений весьма слаб и дает только логарифмическую поправку к $\langle\boldsymbol{R}^2\rangle$.}.

Оценить зависимость $v$ от температуры можно следующим образом. Предположим, что взаимодействие между сегментами описывается потенциальной энергией $u(\boldsymbol{r})$, зависящей только от расстояния между ними. Тогда второй вириальный коэффициент оценивается стандартной формулой для неидеального газа [22]:
\begin{equation}
v = \int\!\mathrm{d}\boldsymbol{r}\left[1 - \exp\left(-\frac{u(\boldsymbol{r})}{k_B T}\right)\right]. \tag{2.92}
\end{equation}
Обычно $u(\boldsymbol{r})$ состоит из потенциала твердого ядра $u_{\mathrm{hard}}(\boldsymbol{r})$ и слабого потенциала притяжения $u_{\mathrm{attr}}(\boldsymbol{r})$ (рис.~2.8). В таком случае $v$ оценивается соотношением
\begin{equation}
v = \int\!\mathrm{d}\boldsymbol{r}\left[1 - \exp\left(-\frac{u_{\mathrm{hard}}(\boldsymbol{r})}{k_B T}\right)\left(1 - \frac{u_{\mathrm{attr}}(\boldsymbol{r})}{k_B T}\right)\right] = A - \frac{B}{T}, \tag{2.93}
\end{equation}
где $A$ и $B$ являются константами, не зависящими от температуры. Для такой модели $\Theta$-температура определяется как $B/A$, а формулу (2.93) можно записать в виде [16]
\begin{equation}
v = v^{(0)}\left(1 - \frac{\Theta}{T}\right). \tag{2.94}
\end{equation}

Теперь если учесть взаимодействие, описываемое формулой (2.87), функция распределения для $\boldsymbol{R}_n$ становится равной
\begin{equation}
\Psi[\boldsymbol{R}_n] \propto \exp\left[-\frac{3}{2b^2}\int_0^N\!\mathrm{d}n\left(\frac{\partial\boldsymbol{R}_n}{\partial n}\right)^2 - \frac{1}{2}v\int_0^N\!\mathrm{d}n\int_0^N\!\mathrm{d}m\,\delta(\boldsymbol{R}_n - \boldsymbol{R}_m)\right]. \tag{2.95}
\end{equation}
Эта модель включает только два параметра: $b$, который представляет взаимодействие между близко расположенными вдоль цепи сегментами, и $v$, описывающий взаимодействие далеко отстоящих друг от друга по цепи сегментов. Основным допущением такой двухпараметрической модели является представление о существовании резкого отличия между указанными мелко- и крупномасштабным взаимодействиями. Хотя справедливость этого допущения и может подвергаться сомнению, как в работе [23], обычно считается, что оно является хорошей отправной точкой для анализа проблемы исключенного объема.

\subsubsection*{2.5.3. ТЕОРЕТИЧЕСКИЕ ПОДХОДЫ}

Обсудим теперь статистические свойства цепи, описываемой формулой (2.95). Ограничимся обсуждением случая $v > 0$, т.~е. случаем, когда преобладает отталкивание. Случаи $v = 0$ и $v < 0$ обсуждаются в работах [24, 25].

\textit{Упрощенная теория.} Оригинальная идея Флори [15] при расчете размеров полимерной цепи состояла в том, чтобы учесть баланс двух эффектов: разбухания полимерной молекулы, обусловленного преобладанием отталкивания при взаимодействиях типа исключенного объема \dots

\textit{Теория среднего поля.} Упрощенная теория, рассмотренная в первой части этого раздела, может быть усовершенствована с использованием расчетов, основанных на понятии среднего поля [41]. Рассмотрим ансамбль цепей, концы которых фиксированы в пространстве: один в начале координат, а другой~--- в точке $\boldsymbol{R}$. Пусть $\bar{c}(\boldsymbol{r})$ есть средняя плотность сегментов в точке $\boldsymbol{r}$:
\begin{equation}
\bar{c}(\boldsymbol{r}) = \int_0^N\!\mathrm{d}n\,\langle\delta(\boldsymbol{r} - \boldsymbol{R}_n)\rangle. \tag{2.114}
\end{equation}
Тогда сегмент в точке $\boldsymbol{r}$ находится под воздействием потенциала среднего поля $\bar{c}(\boldsymbol{r})vk_B T$, а, следовательно, статистическое распределение для рассматриваемой полимерной цепи приобретает вид
\begin{equation}
\Psi_{\mathrm{MF}}[\boldsymbol{R}_n] \propto \exp\left[-\frac{3}{2b^2}\int_0^N\!\mathrm{d}n\left(\frac{\partial\boldsymbol{R}_n}{\partial n}\right)^2 - v\int_0^N\!\mathrm{d}n\,\bar{c}(\boldsymbol{R}_n)\right]. \tag{2.115}
\end{equation}
Для такой функции распределения функция Грина $G(\boldsymbol{R}, 0, n)$ рассчитывается из уравнения (см. уравнение (2.53))
\begin{equation}
\left[\frac{\partial}{\partial n} - \frac{b^2}{6}\frac{\partial^2}{\partial\boldsymbol{R}^2} + v\bar{c}(\boldsymbol{R})\right]G(\boldsymbol{R}, 0, n) = \delta(\boldsymbol{R})\delta(n). \tag{2.116}
\end{equation}
Если $G(\boldsymbol{R}, 0, n)$ найдена, то плотность сегментов $\bar{c}(\boldsymbol{r})$ рассчитывается как
\begin{equation}
\bar{c}(\boldsymbol{r}) = \frac{1}{G(\boldsymbol{R}, 0, N)}\int_0^N\!\mathrm{d}n\,G(\boldsymbol{R}, \boldsymbol{r}, N - n)G(\boldsymbol{r}, 0, n). \tag{2.117}
\end{equation}
Уравнения (2.116) и (2.117) образуют замкнутую систему для нахождения $G(\boldsymbol{R}, 0, N)$. Теория включает в себя достаточно громоздкие расчеты, но на ее основе могут быть предсказаны многие интересные черты поведения цепи с учетом исключенного объема. Показано, например, что структурный фактор $g(k)$ в области высоких значений $k$ имеет вид
\begin{equation}
g(k) \propto k^{-5/3} \quad \text{при } kR_g \gg 1, \tag{2.118}
\end{equation}
что можно получить также на основе модели однородного растяжения. Космас и Фрид [42] дополнили эту модель учетом флуктуаций, но, как оказалось, это не изменяет существенно основных результатов.

\textit{Теория ренормализационной группы.} Теория среднего поля пренебрегает тем обстоятельством, что столкновения между сегментами в полимерной молекуле сильно коррелированы: весьма вероятно, что столкновение любой пары сегментов вызовет другие столкновения, поскольку эти сегменты связаны цепью. Чтобы учесть эту коррелированность, к рассмотрению полимерных задач был применен метод ренормализационной группы, разработанный при изучении критических явлений. Этот метод, первоначально введенный Вильсоном [43], к полимерным проблемам впервые был применен де Женом [44] и де Клуазо [45].

Особенно отчетливо основной результат формулируется Вильсоном и Фишером [46], заметившими, что в пространствах высокой размерности эффект корреляций становится слабым и применимым становится простое разложение по теории возмущений. Например, если проблема исключенного объема рассматривается для полимерной цепи, помещенной в $d$-мерное евклидово пространство, то параметр разложения $z$, равный $N^2 v / (\sqrt{N}b)^d \propto N^{(4-d)/2}$, становится малым для больших $N$, если $d > 4$. (Размерность 4 называют критической размерностью.) Таким образом удается разработать схему разложения, используя в качестве малого параметра $\varepsilon = 4 - d$. Такая схема дает выражение для $\nu$ вида
\begin{equation}
\nu = \frac{1}{2}\left(1 + \frac{1}{8}\varepsilon + \frac{15}{256}\varepsilon^2 \dots\right). \tag{2.119}
\end{equation}
При $d = 3$ эта формула дает $\nu = 0{,}592$, достаточно близкое к более точному результату [36, 38] $\nu = 0{,}588 \pm 0{,}001$. В любом случае отклонения этих показателей от $0{,}6$ весьма малы.

Реальные методы расчета в теории ренормализационной группы весьма сложны и многообразны. Детали можно найти в работах [18, 45, 47, 48].

\section*{2.I.4. ФУНКЦИОНАЛЬНЫЙ ИНТЕГРАЛ}

Если нижний индекс $n$ у $x_n$ рассматривать как непрерывную переменную, то набор точек $(x_1, x_2, \dots, x_N)$ будет представлять собой непрерывную функцию, а интеграл по набору $(x_1, x_2, \dots, x_N)$ сведется к интегрированию по всем возможным функциям; он называется функциональным интегралом и обозначается символом $\delta x_n$, т.~е.
\begin{equation*}
\int\!\prod_n \mathrm{d}x_n \;\xrightarrow[\text{непрерывный предел}]{}\; \int\!\delta x_n.
\end{equation*}
Подынтегральное выражение в функциональном интеграле получается путем предельного перехода от дискретных переменных к непрерывным. Например, в качестве обобщения соотношения (2.I.16) мы имеем
\begin{align}
\int\!\delta x_n \exp\left[-\frac{1}{2}\int_0^N\!\mathrm{d}n\int_0^N\!\mathrm{d}m\,A_{nm}x_n x_m + \int_0^N\!\mathrm{d}n\,\xi_n x_n\right] \propto \nonumber\\
\propto \exp\left[\max\left(-\frac{1}{2}\int_0^N\!\mathrm{d}n\int_0^N\!\mathrm{d}m\,A_{nm}x_n x_m + \int_0^N\!\mathrm{d}n\,\xi_n x_n\right)\right] \tag{2.I.32}
\end{align}
или
\begin{equation}
\left\langle\exp\left[\int_0^N\!\mathrm{d}n\,\xi_n x_n\right]\right\rangle = \exp\left[\frac{1}{2}\int_0^N\!\mathrm{d}n\int_0^N\!\mathrm{d}m\,(A^{-1})_{nm}\xi_n \xi_m\right], \tag{2.I.33}
\end{equation}
где
\begin{equation}
\int_0^N\!\mathrm{d}m\,(A^{-1})_{nm}A_{mj} = \delta(n - j). \tag{2.I.34}
\end{equation}

\section*{ПРИЛОЖЕНИЕ 2.II. ДИФФЕРЕНЦИАЛЬНОЕ УРАВНЕНИЕ ДЛЯ $G(\boldsymbol{R}, \boldsymbol{R}'; N)$}

В случае $U_e = 0$ функция $G(\boldsymbol{R}, \boldsymbol{R}'; N)$ становится распределением вероятностей для вектора между концами и имеет вид
\begin{equation}
G_0(\boldsymbol{R} - \boldsymbol{R}'; N) = \left(\frac{2\pi Nb^2}{3}\right)^{-3/2}\exp\left(-\frac{3(\boldsymbol{R} - \boldsymbol{R}')^2}{2Nb^2}\right)\Theta(N), \tag{2.II.1}
\end{equation}
где
\begin{equation}
\Theta(N) = \begin{cases} 1, & N > 0, \\ 0, & N < 0. \end{cases} \tag{2.II.2}
\end{equation}
Воспользовавшись преобразованием Фурье (см. уравнение (2.25))
\begin{equation}
G_0(\boldsymbol{R} - \boldsymbol{R}'; N) = \int\!\frac{\mathrm{d}\boldsymbol{k}}{(2\pi)^3}\exp(-i\boldsymbol{k}\cdot(\boldsymbol{R} - \boldsymbol{R}'))\exp\left(-\frac{Nb^2 k^2}{6}\right)\Theta(N), \tag{2.II.3}
\end{equation}
можно легко убедиться в том, что
\begin{align}
\frac{\partial}{\partial N}G_0(\boldsymbol{R} - \boldsymbol{R}'; N) &= \int\!\frac{\mathrm{d}\boldsymbol{k}}{(2\pi)^3}\exp(-i\boldsymbol{k}\cdot(\boldsymbol{R} - \boldsymbol{R}'))\left[-\frac{b^2 k^2}{6}\exp\left(-\frac{Nb^2 k^2}{6}\right)\Theta(N) + \delta(N)\right] = \nonumber\\
&= \frac{b^2}{6}\frac{\partial^2}{\partial\boldsymbol{R}^2}G_0(\boldsymbol{R} - \boldsymbol{R}'; N) + \delta(\boldsymbol{R} - \boldsymbol{R}')\delta(N). \tag{2.II.4}
\end{align}
Чтобы вывести уравнение для $G(\boldsymbol{R}, \boldsymbol{R}'; N)$ при $U_e \ne 0$, воспользуемся соотношением (2.49) при малых $\Delta N$:
\begin{equation}
G(\boldsymbol{R}, \boldsymbol{R}'; N + \Delta N) = \int\!\mathrm{d}\boldsymbol{R}''\,G(\boldsymbol{R}, \boldsymbol{R}''; \Delta N)G(\boldsymbol{R}'', \boldsymbol{R}'; N). \tag{2.II.5}
\end{equation}
При $n$, лежащих между $N$ и $N + \Delta N$, $\boldsymbol{R}_n$ не может находиться далеко от $\boldsymbol{R}$. Поэтому, если $U_e(\boldsymbol{r})$ является гладкой функцией $\boldsymbol{r}$, энергия может быть приближенно выражена в виде
\begin{equation}
\int_N^{N+\Delta N}\!\mathrm{d}n\,U_e(\boldsymbol{R}_n) \approx \Delta N\,U_e(\boldsymbol{R}). \tag{2.II.6}
\end{equation}
Тогда $G(\boldsymbol{R}, \boldsymbol{R}''; \Delta N)$ легко найти, используя выражение (2.46),
\begin{equation}
G(\boldsymbol{R}, \boldsymbol{R}''; \Delta N) = \exp\left(-\frac{1}{k_B T}\Delta N\,U_e(\boldsymbol{R})\right)G_0(\boldsymbol{R} - \boldsymbol{R}''; \Delta N). \tag{2.II.7}
\end{equation}
Согласно формулам (2.II.5) и (2.II.7), имеем
\begin{align}
G(\boldsymbol{R}, \boldsymbol{R}'; N + \Delta N) &= \exp\left(-\frac{1}{k_B T}\Delta N\,U_e(\boldsymbol{R})\right)\times \nonumber\\
&\quad \times \int\!\mathrm{d}\boldsymbol{R}''\,G_0(\boldsymbol{R} - \boldsymbol{R}''; \Delta N)G(\boldsymbol{R}'', \boldsymbol{R}'; N). \tag{2.II.8}
\end{align}
При малых $\Delta N$ $G_0(\boldsymbol{R} - \boldsymbol{R}''; \Delta N)$ имеет острый пик при $\boldsymbol{R} = \boldsymbol{R}''$. Поэтому интеграл в формуле (2.II.8) можно оценить, разлагая $G(\boldsymbol{R}'', \boldsymbol{R}'; N)$ в ряд по $\boldsymbol{r} \equiv \boldsymbol{R} - \boldsymbol{R}''$:
\begin{align}
I &= \int\!\mathrm{d}\boldsymbol{R}''\,G_0(\boldsymbol{R} - \boldsymbol{R}''; \Delta N)G(\boldsymbol{R}'', \boldsymbol{R}'; N) = \nonumber\\
&= \int\!\mathrm{d}\boldsymbol{r}\,G_0(\boldsymbol{r}; \Delta N)G(\boldsymbol{R} - \boldsymbol{r}, \boldsymbol{R}'; N) = \nonumber\\
&= \int\!\mathrm{d}\boldsymbol{r}\,G_0(\boldsymbol{r}; \Delta N)\left(1 - r_\alpha\frac{\partial}{\partial R_\alpha} + \frac{1}{2}r_\alpha r_\beta \frac{\partial^2}{\partial R_\alpha \partial R_\beta}\right)G(\boldsymbol{R}, \boldsymbol{R}'; N). \tag{2.II.9}
\end{align}
Так как
\begin{equation}
\int\!\mathrm{d}\boldsymbol{r}\,G_0(\boldsymbol{r}; \Delta N)r_\alpha = 0, \quad \int\!\mathrm{d}\boldsymbol{r}\,G_0(\boldsymbol{r}; \Delta N)r_\alpha r_\beta = \Delta N\frac{b^2}{3}\delta_{\alpha\beta}, \tag{2.II.10}
\end{equation}
то интеграл становится равным
\begin{equation}
I = \left(1 + \frac{1}{6}\Delta Nb^2 \frac{\partial^2}{\partial\boldsymbol{R}^2}\right)G(\boldsymbol{R}, \boldsymbol{R}'; N). \tag{2.II.11}
\end{equation}
Таким образом, с точностью до первого порядка по $\Delta N$ уравнение (2.II.8) запишется как
\begin{align}
\left(1 + \Delta N\frac{\partial}{\partial N}\right)G(\boldsymbol{R}, \boldsymbol{R}'; N) &= \left(1 - \frac{1}{k_B T}\Delta N\,U_e(\boldsymbol{R})\right)\times \nonumber\\
&\quad \times \left(1 + \frac{1}{6}\Delta Nb^2\frac{\partial^2}{\partial\boldsymbol{R}^2}\right)G(\boldsymbol{R}, \boldsymbol{R}'; N). \tag{2.II.12}
\end{align}
Сравнивая члены порядка $\Delta N$, получаем
\begin{equation}
\left(\frac{\partial}{\partial N} - \frac{b^2}{6}\frac{\partial^2}{\partial\boldsymbol{R}^2} + \frac{1}{k_B T}U_e(\boldsymbol{R})\right)G(\boldsymbol{R}, \boldsymbol{R}'; N) = 0. \tag{2.II.13}
\end{equation}
Это уравнение справедливо при $N > 0$. Чтобы учесть сингулярность при $N = 0$, используем выражение (2.II.7) для малых $N$, т.~е.
\begin{equation}
G(\boldsymbol{R}, \boldsymbol{R}'; N) = \exp\left(\frac{-1}{k_B T}N U_e(\boldsymbol{R})\right)G_0(\boldsymbol{R} - \boldsymbol{R}'; N), \tag{2.II.14}
\end{equation}
что и дает $\delta(\boldsymbol{R} - \boldsymbol{R}')\delta(N)$ в правой части формулы (2.53).

\section*{ПРИЛОЖЕНИЕ 2.III. РАСЧЕТ ЭФФЕКТА ИСКЛЮЧЕННОГО ОБЪЕМА ПО ТЕОРИИ ВОЗМУЩЕНИЙ}

Пусть $\langle\dots\rangle_0$ обозначает среднее по функции распределения идеальной цепи; тогда для цепи с исключенным объемом
\begin{equation}
\langle\boldsymbol{R}^2\rangle = \frac{\langle\exp(-U_1)(\boldsymbol{R}_N - \boldsymbol{R}_0)^2\rangle_0}{\langle\exp(-U_1)\rangle_0}, \tag{2.III.1}
\end{equation}
где
\begin{equation}
U_1 = v\int_0^N\!\mathrm{d}m\int_0^m\!\mathrm{d}n\,\delta(\boldsymbol{R}_n - \boldsymbol{R}_m) \tag{2.III.2}
\end{equation}
(здесь $k_B T$ положено равным единице). Сначала рассчитаем
\begin{align}
\langle\boldsymbol{R}^2\rangle &= \frac{\langle(1 - U_1)(\boldsymbol{R}_N - \boldsymbol{R}_0)^2\rangle_0}{\langle 1 - U_1\rangle_0} = \nonumber\\
&= \langle(\boldsymbol{R}_N - \boldsymbol{R}_0)^2\rangle_0(1 + \langle U_1\rangle_0) - \langle U_1(\boldsymbol{R}_N - \boldsymbol{R}_0)^2\rangle_0 = \nonumber\\
&= Nb^2\left(1 + v\int_0^N\!\mathrm{d}m\int_0^m\!\mathrm{d}n\,\langle\delta(\boldsymbol{R}_n - \boldsymbol{R}_m)\rangle_0\right) - \nonumber\\
&\quad - v\int_0^N\!\mathrm{d}m\int_0^m\!\mathrm{d}n\,\langle\delta(\boldsymbol{R}_n - \boldsymbol{R}_m)(\boldsymbol{R}_N - \boldsymbol{R}_0)^2\rangle_0. \tag{2.III.3}
\end{align}
\begin{equation}
I = \langle\delta(\boldsymbol{R}_n - \boldsymbol{R}_m)(\boldsymbol{R}_N - \boldsymbol{R}_0)^2\rangle_0, \tag{2.III.4}
\end{equation}
что с использованием функции Грина $G_0(\boldsymbol{R}, \boldsymbol{R}'; N)$ для свободного пространства (см. уравнения (2.50) и (2.51)) запишется в виде
\begin{align}
I &= \int\!\mathrm{d}\boldsymbol{R}_N\mathrm{d}\boldsymbol{R}_m\mathrm{d}\boldsymbol{R}_n\mathrm{d}\boldsymbol{R}_0\,G_0(\boldsymbol{R}_N - \boldsymbol{R}_m; N - m)\times \nonumber\\
&\quad \times G_0(\boldsymbol{R}_m - \boldsymbol{R}_n; m - n)G_0(\boldsymbol{R}_n - \boldsymbol{R}_0; n)\delta(\boldsymbol{R}_n - \boldsymbol{R}_m)(\boldsymbol{R}_N - \boldsymbol{R}_0)^2 = \nonumber\\
&= \int\!\mathrm{d}\boldsymbol{R}_N\mathrm{d}\boldsymbol{R}_n\mathrm{d}\boldsymbol{R}_0\,G_0(\boldsymbol{R}_N - \boldsymbol{R}_n; N - m)G_0(0; m - n)\times \nonumber\\
&\quad \times G_0(\boldsymbol{R}_n - \boldsymbol{R}_0; n)(\boldsymbol{R}_N - \boldsymbol{R}_0)^2. \tag{2.III.5}
\end{align}
Чтобы взять этот интеграл, запишем $\boldsymbol{R}_N - \boldsymbol{R}_0$ как $(\boldsymbol{R}_N - \boldsymbol{R}_n) - (\boldsymbol{R}_n - \boldsymbol{R}_0)$ и проинтегрируем вначале по $\boldsymbol{R}_N$, а затем по $\boldsymbol{R}_n$. В результате получим
\begin{equation}
I = G_0(0; m - n)(N - m + n)b^2. \tag{2.III.6}
\end{equation}
Следовательно,
\begin{equation}
\langle U_1(\boldsymbol{R}_N - \boldsymbol{R}_0)^2\rangle_0 = v\int_0^N\!\mathrm{d}m\int_0^m\!\mathrm{d}n\,(N - m + n)b^2 G_0(0; m - n). \tag{2.III.7}
\end{equation}
Аналогично
\begin{equation}
\langle U_1\rangle_0 = v\int_0^N\!\mathrm{d}m\int_0^m\!\mathrm{d}n\,G_0(0; m - n). \tag{2.III.8}
\end{equation}
Интегралы в выражениях (2.III.7) и (2.III.8) расходятся, но расходящиеся члены сокращаются друг с другом в формуле для $\langle\boldsymbol{R}^2\rangle$. В конечном счете из соотношений (2.III.3), (2.III.7) и (2.III.8) мы имеем выражение
\begin{align}
\langle\boldsymbol{R}^2\rangle &= Nb^2\left(1 + v\int_0^N\!\mathrm{d}m\int_0^m\!\mathrm{d}n\,G_0(0; m - n)(m - n)\right) = \nonumber\\
&= Nb^2\left(1 + v\int_0^N\!\mathrm{d}m\int_0^m\!\mathrm{d}n\left(\frac{3}{2\pi(m - n)b^2}\right)^{3/2}(m - n)\right), \tag{2.III.9}
\end{align}
которое сходится и дает первый член в формуле (2.100).

Затем выведем формулу (2.108). Интеграл, который мы должны теперь оценить, имеет вид
\begin{equation}
J = \left\langle\left(\frac{\partial\boldsymbol{R}_n}{\partial n}\right)^2(\boldsymbol{R}_N - \boldsymbol{R}_0)^2\right\rangle'. \tag{2.III.10}
\end{equation}
Чтобы вычислить его, заменим $\partial\boldsymbol{R}_n/\partial n$ на $\boldsymbol{R}_n - \boldsymbol{R}_{n-1}$ и перепишем $\boldsymbol{R}_N - \boldsymbol{R}_0$ как
\begin{equation}
\boldsymbol{R}_N - \boldsymbol{R}_0 = (\boldsymbol{R}_N - \boldsymbol{R}_n) + (\boldsymbol{R}_n - \boldsymbol{R}_{n-1}) + (\boldsymbol{R}_{n-1} - \boldsymbol{R}_0). \tag{2.III.11}
\end{equation}
Так как корреляция между $\boldsymbol{R}_N - \boldsymbol{R}_n$, $\boldsymbol{R}_n - \boldsymbol{R}_{n-1}$ и $\boldsymbol{R}_{n-1} - \boldsymbol{R}_0$ отсутствует, выражение (2.III.10) приобретает вид
\begin{align}
J &= \langle(\boldsymbol{R}_n - \boldsymbol{R}_{n-1})^2\rangle'\langle(\boldsymbol{R}_N - \boldsymbol{R}_n)^2\rangle' + \langle(\boldsymbol{R}_n - \boldsymbol{R}_{n-1})^4\rangle' + \nonumber\\
&\quad + \langle(\boldsymbol{R}_n - \boldsymbol{R}_{n-1})^2\rangle'\langle(\boldsymbol{R}_{n-1} - \boldsymbol{R}_0)^2\rangle'. \tag{2.III.12}
\end{align}
Поскольку $\langle(\boldsymbol{R}_m - \boldsymbol{R}_n)^2\rangle' = (m - n)b'^2$ и $\langle(\boldsymbol{R}_m - \boldsymbol{R}_n)^4\rangle' = 5/3(m - n)^2 b'^4$, то уравнение (2.III.12) принимает вид
\begin{equation}
J = Nb'^4 + \frac{2}{3}b'^4, \tag{2.III.13}
\end{equation}
что и дает выражение (2.108).

\end{document}